\section{项目背景与实验设计}\label{sec:project-bg}

\subsection{受测模型与三类基准任务}

为了使实验结论兼具小参数稠密模型与大参数混合专家（MoE）模型的代表性，我们选取了四款主流开源及商业接口模型进行跨任务对照（详细公开规格及来源说明见附录~\ref{app:prompts}）。在本地推理与开源部署端，选用 \texttt{Qwen2.5-7B-Instruct}（7.61B 稠密）与 \texttt{Qwen3-8B}（8.2B 稠密，并在代码反馈补充实验中显式关闭额外思考模式）；在稀疏激活大模型端，选用 \texttt{DeepSeek-V4.1-Flash-1017}（552B 骨干，预填充激活 8B / 生成激活 16B）与 \texttt{HY4 preview}（770B 骨干，每 Token 激活 49B）。

\noindent\begin{minipage}[t]{0.52\linewidth}
\vspace{0pt}
任务设置覆盖课程要求的\textbf{三类典型 AI 交互场景}：第一类为数学推理与学科选择题，采用 \texttt{GSM8K}（小学算术应用题，200 题）、\texttt{MATH500}（竞赛数学题，200 题）与 \texttt{MMLU-Pro}（多学科高难选择题，200 题）；第二类为长上下文文档问答，基于 \texttt{NQ-Open} 构建 200 道问题与 10/20 篇干扰文档组合及闭卷基线；第三类为 Python 代码生成与多轮修复，分别基于 \texttt{CodePromptEval}（清洗审核后保留 194 题）与 \texttt{HumanEval+ Mini}（164 题）开展独立采样与多轮执行反馈对照。
\end{minipage}\hfill
\begin{minipage}[t]{0.45\linewidth}
\vspace{0pt}
\centering
\captionof{table}{受测模型与任务范围概览}\label{tab:models-tasks}
\ReportTableStyle
\begin{tabularx}{\linewidth}{@{}>{\raggedright\arraybackslash}p{1.9cm} >{\raggedright\arraybackslash}X@{}}
\toprule
\textbf{维度} & \textbf{实验范围} \\
\midrule
受测模型 & Qwen2.5 7B、Qwen3 8B、DeepSeek Flash、HY4 preview \\
推理/选择 & GSM8K、MATH500、MMLU-Pro（各 200 题） \\
文档问答 & NQ-Open 200 题（10/20 篇材料 + 闭卷） \\
代码任务 & CodePromptEval 194 题；HumanEval+ 164 题 \\
\bottomrule
\end{tabularx}
\end{minipage}

\subsection{31 种提示策略与五组实验矩阵}

在推理、文档问答与代码生成实验中，我们系统构造了覆盖\textbf{回答组织、语言模板、长度预算与表达格式}四大族系的 \textbf{31 种提示策略}（完整提示文本见附录~\ref{app:prompts} 表~\ref{tab:prompts-31}）。回答组织族包含只给答案（\texttt{NoCoT}）、默认分步（\texttt{DefaultCoT}）、简洁分步（\texttt{BeConcise}）以及步数限制（\texttt{StepLimit(1/2/5)}）；语言模板族包含中文作答（\texttt{ChineseCoT}）及其四档长度组合；长度预算族涵盖词数 7 档（1 至 100 词）、字母数 4 档（10 至 500 字母）和 Token 数 4 档（10 至 500 Token）；表达格式族则包含项目符号、缩写去元音、简化语法、移除空格与纯数字算式。

针对文档上下文位置效应与代码修复闭环机制，我们进一步增设了两项专题实验：其一为\textbf{文档证据位置专题}（3,600 条响应），固定 \texttt{NoCoT} 提示与 20 篇参考文档，在保持问题完全一致的前提下，将唯一包含答案的黄金文档依次置于第 1、3、5、8、10、13、15、18、20 位，精细刻画双模型的位置剖面曲线；其二为\textbf{代码执行反馈补充实验}（24,600 条响应），在共享同一份初始代码草稿的基础上，对比独立重采样、无反馈自行检查、仅返回通过状态、错误定位诊断、单个失败反例与多个失败反例六种处理臂在四轮修改中的演化轨迹。
